% Compile with: pdflatex vector_calculus_study_sheet_ch6_0_to_6_5.tex (run it twice).

\documentclass[12pt, twoside]{article}
\usepackage{amsmath}
\usepackage{amssymb}
\usepackage{geometry}
\usepackage{enumitem}
\usepackage{graphicx}

\geometry{margin=1in}

\makeatletter
\renewcommand{\maketitle}{%
  \begin{center}
    {\LARGE\textbf{\@title}}\\[0.4em]
  \end{center}
  \vspace{-0.5em}
}
\makeatother

\title{Study Sheet: Vector Calculus (Chapter 6.0--6.5)\\[0.3em]
{\large\textit{Calculus Volume 3} (OpenStax)}}

\begin{document}

\maketitle

\noindent\textbf{Scope:} Sections 6.0 through 6.5 only (Introduction, Vector Fields, Line Integrals, Conservative Vector Fields, Green's Theorem, Divergence and Curl).

\bigskip

\section*{1) Core formulas you need}

\subsection*{6.0 Introduction (no heavy formulas)}
\begin{itemize}[leftmargin=*]
  \item Big ideas: vector fields, line integrals, work, circulation, flux, and higher-dimensional FTC ideas.
\end{itemize}

\subsection*{6.1 Vector Fields}
\begin{align*}
\mathbf{F}(x,y) &= \langle P(x,y),Q(x,y)\rangle = P\,\mathbf{i}+Q\,\mathbf{j},\\
\mathbf{F}(x,y,z) &= \langle P(x,y,z),Q(x,y,z),R(x,y,z)\rangle.
\end{align*}
Conservative (gradient) field:
\[
\mathbf{F}=\nabla f\quad\text{for some potential }f.
\]

\subsection*{6.2 Line Integrals}
Curve: $\mathbf{r}(t)=\langle x(t),y(t),z(t)\rangle$, $a\le t\le b$.
\[
\|\mathbf{r}'(t)\|=\sqrt{(x')^2+(y')^2+(z')^2},\qquad ds=\|\mathbf{r}'(t)\|dt.
\]
Scalar line integral:
\[
\int_C f\,ds=\int_a^b f(\mathbf{r}(t))\,\|\mathbf{r}'(t)\|\,dt.
\]
Vector line integral / work:
\[
\int_C \mathbf{F}\cdot d\mathbf{r}=\int_C \mathbf{F}\cdot \mathbf{T}\,ds
=\int_a^b \mathbf{F}(\mathbf{r}(t))\cdot \mathbf{r}'(t)\,dt.
\]
Equivalent differential form:
\[
\int_C \mathbf{F}\cdot d\mathbf{r}=\int_C P\,dx+Q\,dy+R\,dz.
\]
Orientation reversal:
\[
\int_{-C}\mathbf{F}\cdot d\mathbf{r}=-\int_C\mathbf{F}\cdot d\mathbf{r}.
\]
Flux across plane curve $C$:
\[
\int_C \mathbf{F}\cdot \mathbf{N}\,ds
=\int_a^b \mathbf{F}(\mathbf{r}(t))\cdot\mathbf{n}(t)\,dt,
\quad \mathbf{n}(t)=\langle y'(t),-x'(t)\rangle.
\]

\subsection*{6.3 Conservative Vector Fields}
Fundamental Theorem for Line Integrals (FTLI):
\[
\int_C \nabla f\cdot d\mathbf{r}=f(\mathbf{r}(b))-f(\mathbf{r}(a)).
\]
Cross-partial conservative tests on open simply connected domains:
\[
\mathbf{F}=\langle P,Q\rangle\text{ conservative }\Longleftrightarrow P_y=Q_x,
\]
\[
\mathbf{F}=\langle P,Q,R\rangle\text{ conservative }\Longleftrightarrow
P_y=Q_x,\;P_z=R_x,\;Q_z=R_y.
\]

\subsection*{6.4 Green's Theorem}
(Positive orientation = counterclockwise.)
\[
\oint_C P\,dx+Q\,dy=\iint_D\left(Q_x-P_y\right)dA
\quad\text{(circulation/tangential form)}
\]
\[
\oint_C \mathbf{F}\cdot\mathbf{N}\,ds=\iint_D\left(P_x+Q_y\right)dA
\quad\text{(flux/normal form)}
\]
Also: $\oint_C \mathbf{F}\cdot\mathbf{N}\,ds=\oint_C(-Q\,dx+P\,dy)$.

\subsection*{6.5 Divergence and Curl}
Divergence:
\[
\operatorname{div}\mathbf{F}=\nabla\cdot\mathbf{F}
=\begin{cases}
P_x+Q_y,&\mathbf{F}=\langle P,Q\rangle\\
P_x+Q_y+R_z,&\mathbf{F}=\langle P,Q,R\rangle
\end{cases}
\]
Curl:
\[
\operatorname{curl}\mathbf{F}=\nabla\times\mathbf{F}
=\left\langle R_y-Q_z,\;P_z-R_x,\;Q_x-P_y\right\rangle
\quad (\mathbb{R}^3),
\]
\[
\operatorname{curl}\mathbf{F}=(Q_x-P_y)\mathbf{k}
\quad (\mathbb{R}^2).
\]
Useful fact (with standard smoothness + simply connected domain conditions):
\[
\mathbf{F}\text{ conservative }\Longleftrightarrow \operatorname{curl}\mathbf{F}=\mathbf{0}.
\]

\section*{2) Worked examples by problem type}

\begin{enumerate}[leftmargin=*, itemsep=1em]
\item \textbf{Vector field evaluation at a point (6.1)}

Given $\mathbf{F}(x,y)=\langle 2y^2+x-4,\cos x\rangle$, find $\mathbf{F}(0,-1)$.
\[
\mathbf{F}(0,-1)=\langle 2(1)+0-4,\cos 0\rangle=\langle -2,1\rangle.
\]

\item \textbf{Scalar line integral (6.2)}

Evaluate $\int_C (x+y)\,ds$ on $\mathbf{r}(t)=\langle t,t\rangle$, $0\le t\le 1$.
\[
f(\mathbf{r}(t))=2t,\quad \|\mathbf{r}'(t)\|=\sqrt{2}.
\]
\[
\int_C (x+y)\,ds=\int_0^1 2t\sqrt{2}\,dt=\sqrt{2}.
\]

\item \textbf{Vector line integral / work (6.2)}

For $\mathbf{F}=\langle -y,x\rangle$, $\mathbf{r}(t)=\langle\cos t,\sin t\rangle$, $0\le t\le\pi$:
\[
\mathbf{F}(\mathbf{r}(t))=\langle-\sin t,\cos t\rangle,
\quad \mathbf{r}'(t)=\langle-\sin t,\cos t\rangle.
\]
\[
\int_C\mathbf{F}\cdot d\mathbf{r}=\int_0^{\pi}1\,dt=\pi.
\]

\item \textbf{Use FTLI (6.3)}

Evaluate $\int_C \nabla f\cdot d\mathbf{r}$ where $f(x,y)=x^2+2y^2$,
from $(0,0)$ to $(2,2)$.
\[
\int_C\nabla f\cdot d\mathbf{r}=f(2,2)-f(0,0)=(4+8)-0=12.
\]

\item \textbf{Find potential function (6.3)}

For $\mathbf{F}=\langle 2xy^3,\;3x^2y^2+\cos y\rangle$, find $f$ with $\nabla f=\mathbf{F}$.
\[
f_x=2xy^3\Rightarrow f=x^2y^3+h(y).
\]
Then
\[
f_y=3x^2y^2+h'(y)=3x^2y^2+\cos y\Rightarrow h'(y)=\cos y\Rightarrow h(y)=\sin y.
\]
So one potential is
\[
f(x,y)=x^2y^3+\sin y.
\]

\item \textbf{Conservative test with domain check (6.3)}

$\mathbf{F}(x,y)=\left\langle x\ln y,\frac{x^2}{2y}\right\rangle$ on $y>0$.
\[
P_y=\frac{x}{y},\quad Q_x=\frac{x}{y}\Rightarrow P_y=Q_x.
\]
Domain $y>0$ is open and simply connected, so $\mathbf{F}$ is conservative.

\item \textbf{Green's theorem: circulation form (6.4)}

For $\mathbf{F}=\langle -y,x\rangle$ and $C$ the unit circle (CCW):
\[
\oint_C P\,dx+Q\,dy=\iint_D(Q_x-P_y)\,dA
=\iint_D(1-(-1))\,dA=2\,\text{Area}(D)=2\pi.
\]

\item \textbf{Green's theorem: flux form (6.4)}

For $\mathbf{F}=\langle x,y\rangle$ and $C$ circle radius $r$ (CCW):
\[
\oint_C\mathbf{F}\cdot\mathbf{N}\,ds
=\iint_D(P_x+Q_y)\,dA
=\iint_D 2\,dA=2\pi r^2.
\]

\item \textbf{Divergence and curl (6.5)}

Let $\mathbf{F}(x,y,z)=\langle e^x,yz,-yz^2\rangle$.
\[
\operatorname{div}\mathbf{F}=\frac{\partial}{\partial x}(e^x)+\frac{\partial}{\partial y}(yz)+\frac{\partial}{\partial z}(-yz^2)
=e^x+z-2yz.
\]
At $(0,2,-1)$:
\[
\operatorname{div}\mathbf{F}=1-1+4=4.
\]
\end{enumerate}

\section*{3) Practice problems with full solutions}

\subsection*{Problem A (scalar line integral)}
Evaluate $\int_C (x^2+y^2)\,ds$ for $\mathbf{r}(t)=\langle\cos t,\sin t\rangle$, $0\le t\le\pi$.

\textbf{Solution:}
\[
f(\mathbf{r}(t))=\cos^2 t+\sin^2 t=1,\quad \|\mathbf{r}'(t)\|=1.
\]
\[
\int_C (x^2+y^2)\,ds=\int_0^{\pi}1\cdot 1\,dt=\pi.
\]

\subsection*{Problem B (vector line integral/work)}
Evaluate $\int_C \mathbf{F}\cdot d\mathbf{r}$ for $\mathbf{F}=\langle x,y\rangle$ and
$\mathbf{r}(t)=\langle t,t^2\rangle$, $0\le t\le 1$.

\textbf{Solution:}
\[
\mathbf{F}(\mathbf{r}(t))=\langle t,t^2\rangle,
\quad \mathbf{r}'(t)=\langle 1,2t\rangle.
\]
\[
\mathbf{F}(\mathbf{r}(t))\cdot\mathbf{r}'(t)=t+2t^3.
\]
\[
\int_C\mathbf{F}\cdot d\mathbf{r}=\int_0^1 (t+2t^3)dt
=\left[\frac{t^2}{2}+\frac{t^4}{2}\right]_0^1=1.
\]

\subsection*{Problem C (FTLI)}
Let $f(x,y,z)=x^2y+z^2$. Evaluate $\int_C\nabla f\cdot d\mathbf{r}$ from $(1,1,0)$ to $(2,3,1)$.

\textbf{Solution:}
\[
\int_C\nabla f\cdot d\mathbf{r}=f(2,3,1)-f(1,1,0)
=(4\cdot 3+1)-(1\cdot 1+0)=13-1=12.
\]

\subsection*{Problem D (conservative test in 2D)}
Determine whether $\mathbf{F}(x,y)=\langle 2xy, x^2+1\rangle$ is conservative on $\mathbb{R}^2$.

\textbf{Solution:}
\[
P_y=2x,\qquad Q_x=2x \Rightarrow P_y=Q_x.
\]
Domain $\mathbb{R}^2$ is open and simply connected, so $\mathbf{F}$ is conservative.
A potential function:
\[
f_x=2xy\Rightarrow f=x^2y+h(y),\quad f_y=x^2+h'(y)=x^2+1\Rightarrow h'(y)=1.
\]
So one choice is $f(x,y)=x^2y+y$.

\subsection*{Problem E (Green circulation)}
Compute $\oint_C(-y\,dx+x\,dy)$ where $C$ is the triangle with vertices $(0,0),(1,0),(0,1)$ oriented counterclockwise.

\textbf{Solution:}
\[
P=-y,\quad Q=x\Rightarrow Q_x-P_y=1-(-1)=2.
\]
Region area is $\frac12$, so
\[
\oint_C P\,dx+Q\,dy=\iint_D 2\,dA=2\cdot\frac12=1.
\]

\subsection*{Problem F (Green flux)}
Compute $\oint_C\mathbf{F}\cdot\mathbf{N}\,ds$ for $\mathbf{F}=\langle x,2y\rangle$ on the circle $x^2+y^2=4$ (CCW).

\textbf{Solution:}
\[
P_x+Q_y=1+2=3,
\quad \text{Area}(D)=\pi(2)^2=4\pi.
\]
\[
\oint_C\mathbf{F}\cdot\mathbf{N}\,ds=\iint_D 3\,dA=3(4\pi)=12\pi.
\]

\subsection*{Problem G (divergence + curl in 2D)}
For $\mathbf{F}(x,y)=\langle x^2y,\,y-x y^2\rangle$, find $\operatorname{div}\mathbf{F}$ and $\operatorname{curl}\mathbf{F}$.

\textbf{Solution:}
\[
\operatorname{div}\mathbf{F}=P_x+Q_y=2xy+(1-2xy)=1.
\]
\[
\operatorname{curl}\mathbf{F}=(Q_x-P_y)\mathbf{k}=(-y^2-x^2)\mathbf{k}.
\]

\section*{4) Step-by-step workflow for vector calculus problems (through 6.5)}

\begin{enumerate}[leftmargin=*, itemsep=0.6em]
  \item \textbf{Classify the problem first:}
  vector field eval/sketch, scalar line integral, vector line integral/work, conservative test, FTLI, Green circulation/flux, divergence/curl.

  \item \textbf{Set up geometry and orientation:}
  identify curve/surface region, parameterization, endpoints, and orientation (especially for vector line integrals and Green's theorem).

  \item \textbf{Choose the fastest theorem:}
  \begin{itemize}[leftmargin=*]
    \item If integrand is $\nabla f\cdot d\mathbf{r}$, use FTLI.
    \item If closed plane curve and suitable region: try Green's theorem.
    \item If simple direct parameterization is clean: integrate in $t$.
  \end{itemize}

  \item \textbf{Compute core derivatives carefully:}
  $\mathbf{r}'(t)$, $\|\mathbf{r}'(t)\|$, partials $P_x,P_y,Q_x,Q_y,\dots$.

  \item \textbf{For conservative checks:}
  verify cross-partials \emph{and} check domain is open/simply connected before concluding conservative.

  \item \textbf{For line integrals:}
  convert everything to one variable $t$ and integrate on correct bounds.

  \item \textbf{For Green's theorem:}
  use counterclockwise orientation convention; if clockwise, insert a minus sign.

  \item \textbf{Interpret sign/units:}
  negative work, positive/negative flux, zero circulation, etc.

  \item \textbf{Sanity check:}
  compare with expected symmetry/physical behavior (e.g., conservative closed-loop circulation should be 0).
\end{enumerate}

\end{document}
